\begin{codebox}
\codeline{\textcolor{cmtgray}{\#\ ==============================}}
\codeline{\textcolor{cmtgray}{\#\ 4.\ 知识熵减的视角}}
\codeline{\textcolor{cmtgray}{\#\ ==============================}}
\codeline{\textcolor{cmtgray}{\#\ 企业知识的自然状态是高熵的：}}
\codeline{\textcolor{cmtgray}{\#\ \ \ 分散在文档、数据库、员工头脑中}}
\codeline{\textcolor{cmtgray}{\#\ \ \ 同一概念多种叫法，同一名称多种含义}}
\codeline{\textcolor{cmtgray}{\#\ \ \ 规则隐含在经验里，无法被系统校验}}
\codeblank{}
\codeline{\textcolor{cmtgray}{\#\ 本体工程的本质是知识熵减：}}
\codeline{\textcolor{cmtgray}{\#\ \ \ 隐性\ \(\rightarrow\)\ 显式：把头脑中的概念写成形式化定义}}
\codeline{\textcolor{cmtgray}{\#\ \ \ 分散\ \(\rightarrow\)\ 统一：建立全企业共享的概念词汇表}}
\codeline{\textcolor{cmtgray}{\#\ \ \ 模糊\ \(\rightarrow\)\ 精确：用公理排除非法状态}}
\codeline{\textcolor{cmtgray}{\#\ \ \ 静态\ \(\rightarrow\)\ 可推理：让机器从已有知识推导新知识}}
\codeblank{}
\codeline{\textcolor{cmtgray}{\#\ 熵减的产出：}}
\codeline{\textcolor{cmtgray}{\#\ \ \ 形式化本体（TBox）+\ 实例数据（ABox）}}
\codeline{\textcolor{cmtgray}{\#\ \ \ =\ 可查询、可推理、可复用、可解释的知识结构}}
\codeblank{}
\end{codebox}
